\section{Causal simulation and a joint attained-error law}
\begin{proposition}[Typed common-task composition]\label{prop:simulation}
A simulator from a source experiment to a target experiment must specify a parameter-independent strategy lift, report transformation, internal state, clock, task readout and physical-command compatibility. Suppose it uses at most $S$ retained states, a target observer uses $M$, and their entire stopped scored transcript laws differ by at most $\eps$ in total variation, with convention $\|P-Q\|_{\TV}=\sup_A|P(A)-Q(A)|$. For loss in $[0,L]$, the composed observer uses at most $SM$ states and its risk differs by at most $L\eps$. Its expected cost bound requires its own additive-cost or killed-potential certificate. Two compatible causal simulators have total deficiency at most the sum of their deficiencies and a product state upper bound.
\end{proposition}
\begin{proof}
The simulator and observer retain their ordered pair of states; clocks and any continuing random seeds are part of those states. Compatibility makes the composition a legal causal procedure. For a common bounded scored loss, the layer-cake formula and the total-variation convention give the risk bound. Contractivity of total variation under a common Markov kernel and the triangle inequality give the two-simulator bound on complete stopped scored laws. Nonnegative costs must instead be summed along actual executed calls or bounded by separately certified stopped potentials. A small terminal-law discrepancy alone can coexist with arbitrarily large expected waiting. Reverse risk bounds need a reverse simulator or a separate attained lower construction.
\end{proof}

The state product is an implementation upper, not a claim of optimal simulator size. Root-mean-square readout errors add by the triangle inequality before squaring; deterministic errors may alternatively use the bounded-loss estimate following Proposition~\ref{prop:relative}. Neither convention makes an allowed upper bound on calibration or simulation error an unavoidable risk lower.

\begin{theorem}[Same-task attained calibration and erasure]\label{thm:joint}
Use the exact balanced class $D_0=1$, $B=2$. Add an initial paid observation of an independent fair sign $E_0$ erased with probability $2\eps$, where $0\le\eps\le1/2$, and a final paid audit of
\[
 (Z,E_0,s\delta),\qquad s\in\{-1,1\},\quad\delta\ge0.
\]
The fixed unknown sign $s$ affects no pre-audit kernel. Conditional on the initial reported tag, the entire core experiment contains no additional information about $E_0$ on the erased branch; in particular its physical kernels may not read that hidden sign. All output coordinates are committed before this audit. Use squared Euclidean loss on the product carrier. For $N\ge\kappa+2$ and $M\ge13$, put
\[
 b=(N-2)/\kappa,\qquad m=\lfloor(M-1)/3\rfloor-2.
\]
Then
\begin{equation}\label{eq:joint-bounds}
 \Qm_M(b)+2\eps+\delta^2
 \le R_{\rm joint}^\Theta(N,M)
 \le\Qm_m(b)+2\eps+\delta^2.
\end{equation}
If \eqref{eq:rectangle} holds, this is a matching law
\begin{equation}\label{eq:joint-rate}
 R_{\rm joint}^\Theta(N,M)
 \asymp F((N-2)/\kappa,M)+\delta^2+\eps,
\end{equation}
with the geometry constants specified earlier. The early report channel has causal deficiency \emph{exactly} $\eps$ relative to the perfect-sign channel when its virtual sign must be delivered before any core call.
\end{theorem}
\begin{proof}
For the upper, use one initial state and three tagged copies of the $2+m$-state core observer: tags $+$, $-$ and erasure. Their number is $1+3(2+m)\le M$. Output the reported sign on known branches, zero on erasure, and zero for the inaccessible calibration sign. The balanced core upper has the same risk after every preceding history, hence its contribution is $\Qm_m(b)$. The other contributions are exactly $2\eps$ and $\delta^2$. Initial observation and final audit each cost one.

For the lower, use the fixed reference core. Expected core cost is at most $N-2$, and its terminal projected dictionary has at most $M$ values, so Lemma~\ref{lem:lower} gives $\Qm_M(b)$. On the erased branch, induction over the causal core kernels and independent observer coins leaves $E_0$ fair conditional on all pre-audit information. Any sign estimate therefore has conditional squared loss at least one. Finally average the two indistinguishable alternatives $s=\pm1$: for any committed last coordinate $r$, the average loss is $\delta^2+r^2\ge\delta^2$. The pre-audit law, core lower and erased-sign lower are identical in these two alternatives. Their three lower contributions therefore hold simultaneously in one worst-case total-risk bound, not for incompatible separately optimized tasks. For $M\ge13$, $m\ge M/8$; Proposition~\ref{prop:geometry} proves \eqref{eq:joint-rate}.

On an erasure, any immediately supplied virtual sign is independent of the true fair sign. The perfect target has event \{virtual sign equals $E_0$\} with probability one, whereas every source simulator has this probability at most $1-\eps$. This gives deficiency at least $\eps$ for the complete scored transcript. Reporting the known sign when available and an independent fair sign otherwise disagrees with the perfect report with probability $\eps$. A coupling, followed by the same causal core strategy and task audit, gives the opposite inequality. The early completion requirement forbids waiting for the final audit to learn $E_0$.
\end{proof}

The joint theorem is an attained construction for one specified scored task, not a universal additive theorem for arbitrary error tolerances. Numerics adds its own relative-acceptance and readout certificates to this construction. Persistent labels, the $S$ simulator states, scratch, program description, pulse precision, raw acquisitions and physical time remain distinct accounts.
